% !TEX TS-program = lualatexmk
% !TEX parameter =  --shell-escape

\documentclass[vespers_book_la_eng.tex]{subfiles}

\ifcsname preamble@file\endcsname
  \setcounter{page}{\getpagerefnumber{M-vespers_book_la_eng_sunday_psalms}}
\fi

\begin{document}
 
 \twosided[pc] %% to allow for the left-right switching of columns needed for bilingual liturgical typesetting (like in missals); it's unclear why this goes BEFORE everything else and not when activating the environment. %% It's OK for paper booklets to not switch, but it's probably best practice to switch.
 
\thispagestyle{empty}

\chapter*{SUNDAYS AT VESPERS.}
\phantomsection
\addcontentsline{toc}{chapter}{Sundays at Vespers.}
 \header{sundays at vespers.}
 
 \section[Throughout the Year.]{THROUGHOUT THE YEAR.}
 
\firstantiphon{7. c2}
\initialscore{1_Dixit_Dominus}

\footnotetext[1]{The words are not repeated, which is always done when the antiphon begins with the first words of the psalm, which continues therefore after the word with which the antiphon ends, whether it was sung in full or simply begun. Regardless, if the words are the same, the antiphon is continued with the psalm.}

\rubrique{\begin{enpars}
Sit at the word ``Sede.'' The cantor sings odd verses; all sing even verses, and this for the hymn and Magnificat as well.
\end{enpars}}

\sdpsalmusbilingue{109}{109_7}{109_en}

\smallscore{an--dixit_dominus_domino--solesmes_1961}

\antiphona{english}{The Lord said to my Lord: Sit thou at my right hand}

\rubrique{\begin{enpars}
Clerics in descending order, or a cantor, intone the subsequent antiphons. The cantor then continues the psalm.
\end{enpars}}

\gscore[2. Ant.]{3. b}{2_Magna_Opera}

\psalmus{110}{110_3b}{110_3a_b}{110_en}

\smallscore{an--magna_opera_domini--solesmes_1961}

\antiphona{english}{Great are the works of the Lord, sought out according to all his wills}

\gscore[3. Ant.]{4. g}{3_Qui_timet_Dominum}

\psalmus{111}{111_4g}{111_4g}{111_en}

\smallscore{an--qui_timet_dominum--solesmes_1961}

\antiphona{english}{He who feareth the Lord shall delight exceedingly in his commandments}

\gscore[4. Ant.]{7. c}{4_Sit_nomen_Domini}

\psalmus{112}{112_7c}{112_7}{112_en}

\smallscore{an--sit_nomen_domini..._in--solesmes_1961}

\antiphona{english}{Blessed be the name of the Lord forever}

\gscore[5. Ant.]{T. pereg.}{5_Deus_autem_noster}

\psalmus{113}{113_per}{113_per}{113_en}

\smallscore{an--deus_autem_noster--solesmes_1961}

\antiphona{english}{But our God is in heaven, he hath done all things whatsoever he would}

\rubrique{\begin{enpars}
Having sung the last antiphon in full, stand when the celebrant does.
\end{enpars}}

\rubrique{\begin{enpars}
The chapter, hymn, and verse are always the same from the IV Sunday after Pentecost until the last Sunday before Advent inclusive, as well as from the II Sunday after the Epiphany until the last before Septuagesima inclusive.
\end{enpars}}

\smalltitle{Chapter.}
\scripture{2 Cor. 1, 3–4.}

\textes{capitulum_pa_la_en}

%% The usual tone and the 3rd tone, used ad libitum, are omitted as we do not use these tones, but they are included in the tex file and the gabc files are available, for the sake of completeness.

\hymnlateng

\label{hy_Lucis_Creator_optime} \phantomsection

\gscore[]{8.}{hy--Lucis_Creator_optime}

\textes{lucis_creator_en}

\rubrique{\begin{enpars}
If a commemoration is made of the Blessed Virgin Mary, the doxology is as follows:
\end{enpars}}

\smallscore{hy--Lucis_Creator_BMV}

\begin{enpars}Jesus, to thee be Glory, who wast born of the Virgin, With the Father and the loving Spirit, in the eternal age. Amen.
\end{enpars}

\smalltitle{2. Another tone ad lib.}

\gscore[]{8.}{hy--Lucis_Creator_optime_2_ad_lib}

\rubrique{Or:}

\smallscore{hy--Lucis_Creator_Optime_2_ad_lib_BMV}

\smalltitle{3. Another tone ad lib.}

\gscore[]{1.}{hy--Lucis_Creator_optime_mode_1}

\rubrique{Or:}

\smallscore{hy--Lucis_Creator_Optime_mode_1_BMV}

\textes{cantor_rubric}
\label{dirigatur}
\smallscore{versiculus_dom_pa}
\begin{enpars} \noindent \vv Direct, O Lord, my prayer.

\noindent \rr As incense in thy sight.
 \end{enpars}
 
 \rubrique{\begin{enpars}
The celebrant intones the \normaltext{Magnificat} antiphon; this is always proper to the day. The cantor continues the \normaltext{Ma\-gnificat.}\end{enpars}}

 \rubrique{\begin{enpars}
The office continues with the \normaltext{Magnificat} antiphon and collect, found under the Sunday or feast, and with the rest of the office as in the Ordinary,
\normaltext{\pageref{M-dv}.}\end{enpars}}

 \section[In Paschal Time.]{IN PASCHAL TIME.}
  \header{sundays at vespers in paschal time.}
  \label{pss_tp} \phantomsection
  
  \smalltitle{Antiphon. 7. c2.}
\gscore[]{}{alleluia_sun_tp_vespers}

\rubrique{This antiphon is sung over all of the psalms. On the I Sunday after Easter, the office is of double major rite, so the antiphon \normaltext{Allelúia, \pageref{duplex_antiphon}} is sung before and after.}

%\rubrique{\begin{otherlanguage}{english}
%Sit at the word ``Sede.'' The cantor sings odd verses; all sing even verses, and this for the hymn and Magnificat as well.
%\end{otherlanguage}}

\psalmus{109}{109_7c2}{109_7}{109_en}

\rubrique{\begin{otherlanguage}{english}
The cantors then begin the psalms at the conclusion of the doxology.
\end{otherlanguage}}

\psalmus{110}{110_7c2}{110_7}{110_en}

\psalmus{111}{111_7c2}{111_7}{111_en}

\psalmus{112}{112_7c2}{112_7}{112_en}

\psalmus{113}{113_7c2}{113_7}{113_en}

\label{duplex_antiphon}

\smallscore{an_alleluia_sun_tp_vespers_solesmes_1961}
\antiphona{english}{Alleluia, alleluia, alleluia}

\end{document}